一条训练曲线摆在你面前：loss 从 2.3 掉到 0.4，中途反弹过两次。这条曲线说明了什么？

如果这个配置你只跑过一次，能说的其实不多。你看到的是一次实现的完整历史；你想说的那句话——``这个配置会收敛到 0.4 附近''——却是关于所有可能历史的断言。两者之间隔着一整个概率模型，而随机过程正是用来填这个空的。

概率论前面几章没有这个麻烦。一个随机变量、或者有限个随机变量组成的向量，谈期望谈分布谈相关，指的始终是同一件事。给不确定性加上时间之后，对象变成一整族随机变量

\[
\set{X_t : t\in T},
\]

换一个等价的说法：样本空间上的一个随机\emph{函数}。它的取值不再是一个数，而是一整条曲线。

\begin{center}
\begin{tikzpicture}[x=1cm,y=1cm,font=\small,>=stealth]
  \draw[black!50,fill=blue!6] (0,0) ellipse (1.15 and 0.85);
  \node[black!70,font=\footnotesize] at (0,1.15) {\(\Omega\)};
  \fill (-0.15,0.05) circle (1.6pt);
  \node[below,font=\footnotesize] at (-0.15,-0.05) {\(\omega\)};
  \draw[->,black!60] (1.35,0.05) -- (2.55,0.05);
  \node[above,font=\footnotesize,black!70] at (1.95,0.1) {\(X\)};
  \draw[black!55] (2.8,0.05) -- (5.0,0.05);
  \fill[blue!70] (3.9,0.05) circle (1.8pt);
  \node[below,font=\footnotesize] at (3.9,-0.05) {\(X(\omega)\)};
  \node[font=\footnotesize,align=center,black!75] at (2.5,-1.15) {随机变量：\\一次实验交付一个数};
  \draw[black!25] (6.1,-1.7) -- (6.1,1.5);
  \begin{scope}[shift={(7.0,0)}]
    \draw[black!50,fill=blue!6] (0,0) ellipse (1.15 and 0.85);
    \node[black!70,font=\footnotesize] at (0,1.15) {\(\Omega\)};
    \fill (-0.15,0.05) circle (1.6pt);
    \node[below,font=\footnotesize] at (-0.15,-0.05) {\(\omega\)};
    \draw[->,black!60] (1.35,0.05) -- (2.55,0.05);
    \node[above,font=\footnotesize,black!70] at (1.95,0.1) {\(X\)};
    \draw[black!55] (2.8,-0.75) -- (5.2,-0.75);
    \draw[black!55] (2.8,-0.75) -- (2.8,0.95);
    \draw[blue!70,thick] (2.9,-0.15) -- (3.2,0.25) -- (3.5,0.0) -- (3.8,0.45) -- (4.1,0.15)
      -- (4.4,0.55) -- (4.7,0.3) -- (5.0,0.6);
    \node[right,font=\footnotesize] at (5.25,-0.75) {\(t\)};
    \node[font=\footnotesize,align=center,black!75] at (3.9,-1.55) {随机过程：\\一次实验交付一整条曲线};
  \end{scope}
\end{tikzpicture}

\emph{同一个随机结果 \(\omega\)，在左边被映成数轴上的一个点，在右边被映成一条完整的函数。随机过程的取值空间不是数轴，而是由 \(T\) 到 \(S\) 的函数组成的轨道空间。}
\end{center}

这张图立刻带来本章反复要用的一组区分。固定 \(\omega\)，你读到的是一条轨道，也就是某一次实验从头到尾的完整历史；固定 \(t\)，你读到的是一个截面，也就是所有可能世界在这一刻的分布。入门阶段的事故大多是这两者串了线：把一次运行的起伏当成总体的方差，把截面分布的收敛当成轨道的收敛，或者反过来用一条曲线的平均值去代表整个分布。后一件事有时是合法的，但它需要一个额外的、经常没人检查的条件。

本章三篇沿着这条线推进：

\begin{itemize}
\tightlist
\item \hyperref[sec:06-Random-Processes/01-Foundations/01]{``随机过程、时间指标与样本路径''} 建立轨道与截面这两个方向，并说明为什么一列边缘分布不足以确定一个过程；
\item \hyperref[sec:06-Random-Processes/01-Foundations/02]{``有限维分布怎样描述一个过程''} 给出补齐依赖结构的标准办法，以及 Kolmogorov 定理保证了什么、没保证什么；
\item \hyperref[sec:06-Random-Processes/01-Foundations/03]{``均值、相关与平稳性''} 退到能从数据里估出来的二阶统计量，并把平稳与遍历这两件常被混为一谈的事分开。
\end{itemize}

三篇都是主线，第三篇的后半段——遍历性——是全章认知落差最大的地方：平稳说的是截面分布不随时间变化，遍历说的是一条足够长的轨道能替代所有轨道的平均。后者才是``跑一次长实验就敢报告期望''的合法性来源，也是 MCMC 之所以能用的理由。

读完本章，你会有一套描述随机过程的通用词汇，但还没有任何一个具体模型。紧接着的\hyperref[ch:047]{``计数过程：事件何时到来、到来了多少''}会把这套词汇用在第一类完整模型上：事件一个接一个到达，问的是何时到、来了多少。
